% Learning Mocha — CS426 final report: shared preamble.
\usepackage[T1]{fontenc}
\usepackage[utf8]{inputenc}
\usepackage[a4paper,margin=2.4cm,headheight=14pt]{geometry}
\usepackage{lmodern}
\usepackage{microtype}
\usepackage{booktabs}
\usepackage{longtable}
\usepackage{tabularx}
\usepackage{array}
\usepackage{xcolor}
\usepackage{graphicx}
\usepackage{listings}
\usepackage{fancyhdr}
\usepackage{titlesec}
\usepackage{enumitem}
\usepackage[breakable,skins]{tcolorbox}
\usepackage{caption}
\usepackage{float}
\usepackage{tikz}
\usetikzlibrary{positioning,arrows.meta,shapes.geometric,fit,calc,backgrounds}
\usepackage[hidelinks]{hyperref}

% --- Mocha palette (mirrors app/src/main/res/values/colors.xml) ---
\definecolor{mochaBrown}{HTML}{6B4F3A}
\definecolor{mochaDeep}{HTML}{4A3428}
\definecolor{mochaCream}{HTML}{FFF7F0}
\definecolor{mochaCreamDark}{HTML}{F3E6DA}
\definecolor{mochaSage}{HTML}{7A8F7A}
\definecolor{mochaText}{HTML}{2C241C}
\definecolor{mochaMuted}{HTML}{6E6258}
\definecolor{mochaError}{HTML}{B54A3C}

\hypersetup{colorlinks=false, pdftitle={Learning Mocha — CS426 Final Report}}

\titleformat{\section}{\Large\bfseries\color{mochaDeep}}{\thesection}{0.7em}{}
\titleformat{\subsection}{\large\bfseries\color{mochaBrown}}{\thesubsection}{0.6em}{}
\titleformat{\subsubsection}{\normalsize\bfseries\color{mochaBrown}}{\thesubsubsection}{0.6em}{}
\titlespacing*{\section}{0pt}{1.6ex plus .4ex}{1.0ex}

\pagestyle{fancy}
\fancyhf{}
\fancyhead[L]{\small\color{mochaMuted}Learning Mocha}
\fancyhead[R]{\small\color{mochaMuted}CS426 Final Project}
\fancyfoot[C]{\small\color{mochaMuted}\thepage}
\renewcommand{\headrulewidth}{0.4pt}
\renewcommand{\footrulewidth}{0pt}

\setlist[itemize]{leftmargin=1.3em,itemsep=1pt,topsep=2pt}
\setlist[enumerate]{leftmargin=1.5em,itemsep=1pt,topsep=2pt}

\lstdefinestyle{mocha}{
  basicstyle=\ttfamily\footnotesize,
  keywordstyle=\color{mochaBrown}\bfseries,
  stringstyle=\color{mochaSage},
  commentstyle=\color{mochaMuted}\itshape,
  numberstyle=\tiny\color{mochaMuted},
  backgroundcolor=\color{mochaCreamDark!45},
  frame=single, rulecolor=\color{mochaCreamDark},
  breaklines=true, breakatwhitespace=false,
  showstringspaces=false, tabsize=2,
  columns=fullflexible, keepspaces=true,
  xleftmargin=0.6em, xrightmargin=0.2em,
  aboveskip=0.8em, belowskip=0.8em,
}
\lstset{style=mocha}

\newtcolorbox{mochabox}[1][]{
  breakable, enhanced, colback=mochaCream, colframe=mochaBrown,
  boxrule=0.5pt, arc=2pt, left=6pt, right=6pt, top=5pt, bottom=5pt,
  fonttitle=\bfseries, #1
}

\newcommand{\code}[1]{\texttt{\small #1}}
\newcolumntype{L}[1]{>{\raggedright\arraybackslash}p{#1}}

% Table-header label used inside TikZ schema nodes.
\newcommand{\thdr}[1]{{\bfseries\color{mochaDeep}#1}}
